\subsection{Dataset}

We evaluate our methods on the MSLesSeg dataset \cite{guarnera_mslesseg_2025},
which contains longitudinal brain MRI scans of multiple sclerosis patients.
The dataset provides MRI images with manual lesion annotations for the
training set and unlabeled scans for the official test set. In total, the dataset includes 75 subjects. The labeled portion consists of 53 patients, each having between one and four longitudinal timepoints. Since the official test set does not provide segmentation labels, we use only the
labeled portion of the dataset for our experiments. To avoid data leakage
across longitudinal scans, the data partitioning is performed at the
\textit{patient level}. We adopt a 5-fold cross-validation strategy, where the 53 patients are
randomly divided into five disjoint folds. In each fold, approximately
80\% of the patients are used for training and the remaining 20\% for
validation. This process is repeated across all five folds such
that each patient appears in the test set exactly once. For evaluation, predictions are generated for each fold independently. Performance metrics are computed per fold and then averaged across all
five folds to obtain the final reported results. Each MRI scan includes three modalities: T1-weighted (T1-w), T2-weighted (T2-w), and FLAIR images. The dataset provided by the MSLesSeg organizers
is already skull-stripped and co-registered to the MNI152 template. All images are resampled to an isotropic voxel spacing of $(1.0, 1.0, 1.0)\,\mathrm{mm}$.

\subsection{Implementation Details}

All models are trained using the \texttt{nnUNet} framework (version 2.1.1).
Training is performed for 150 epochs with a batch size of 2 using the Adam
optimizer and an initial learning rate of 0.01. The default nnUNet configuration uses a combination of Dice loss and cross-entropy loss. Based on this baseline, we implement three variants of
loss functions designed to improve lesion detection:

\begin{itemize}
    \item \textbf{CAT}: Component-Adaptive Tversky term that emphasizes
    small connected components.
    \item \textbf{MIL}: A lesion-level multiple-instance learning term that
    encourages the detection of individual lesions.
    \item \textbf{CATMIL}: A combination of CAT and MIL to jointly improve
    voxel-level segmentation and lesion-level detection.
\end{itemize}

All experiments are conducted on a workstation equipped with an NVIDIA
Quadro RTX 8000 GPU.

\subsection{Compared Methods}

% We compare the proposed loss formulations with several widely used
% 3D medical image segmentation architectures, including nnUNet
% \cite{isensee_nnu-net_2021}}, SegResNet \cite{myronenko_3d_2018}, UNETR \cite{hatamizadeh_unetr_2021}, SwinUNETR \cite{hatamizadeh_swin_2022},
% UMambaEnc \cite{ma_u-mamba_2024}, and UMambaBot \cite{ma_u-mamba_2024}.
% For a fair comparison, all models are trained using the same dataset split
% and experimental protocol.

We compare the proposed loss formulation against several widely used
segmentation loss functions using a U-Net architecture implemented
within the nnUNet framework~\cite{isensee_nnu-net_2021}. In
particular, we benchmark our method against standard losses commonly
used in medical image segmentation, including Dice + Cross Entropy
(nnUNet default), as well as Tversky and Focal Tversky losses,
which are designed to address class imbalance by reweighting false
positives and false negatives. All models are trained using the same
dataset split, preprocessing pipeline, and training protocol to ensure
a fair and consistent comparison. Since the architecture is fixed across all experiments, we denote each
variant according to the loss function used, namely
nnUNet-DiceCE, nnUNet-Tversky,
nnUNet-FocalTversky, and nnUNet-CATMIL (ours).

\subsection{Evaluation Metrics}

To comprehensively evaluate segmentation performance, all methods are
assessed using a combination of voxel-level and lesion-level metrics.
While Dice and HD95 are standard in medical image segmentation, they
primarily reflect global overlap and boundary accuracy, and may not
adequately capture the detection of small and sparse lesions. In
particular, small lesions contribute minimally to voxel-level overlap and
can be overlooked without significantly affecting Dice scores.
Therefore, additional lesion-level and size-aware metrics are included
to provide a more complete evaluation of model behavior, especially in
the sparse lesion regime.

\paragraph{Voxel-level Metrics.}
Voxel-level metrics measure the agreement between predicted and
ground-truth segmentation masks at the voxel level. Let $\Omega$ denote
the voxel domain, $P \subset \Omega$ the set of predicted lesion voxels,
and $G \subset \Omega$ the set of ground-truth lesion voxels.

\begin{itemize}
    \item \textbf{Dice Similarity Coefficient (Dice).} Measures the
    spatial overlap between the predicted segmentation and the
    ground-truth annotation:
    \begin{equation}
    \text{Dice} =
    \frac{2|P \cap G|}{|P| + |G|}.
    \end{equation}

    \item \textbf{Hausdorff Distance (HD95).} Measures the boundary
    discrepancy between prediction and ground truth using the 95th
    percentile of surface distances:
    \begin{equation}
    d(S_P,S_G) =
    \left\{
    \min_{g \in S_G} d(p,g) : p \in S_P
    \right\},
    \end{equation}
    \begin{equation}
    HD_{95}(P,G) =
    \max \left\{
    \mathrm{percentile}_{95}\big(d(S_P,S_G)\big),
    \mathrm{percentile}_{95}\big(d(S_G,S_P)\big)
    \right\},
    \end{equation}
    where $S_P$ and $S_G$ denote the sets of surface voxels of the
    predicted and ground-truth masks, respectively, and $d(\cdot,\cdot)$
    denotes the Euclidean distance.
\end{itemize}

\paragraph{Lesion-level and Error Metrics.}
Voxel-level metrics may not fully reflect lesion detection performance,
especially for small and sparse lesions. Therefore, lesion-level
evaluation is additionally performed using connected component analysis,
where each lesion is defined as a 3D connected component of the binary
mask. Let $\mathcal{G} = \{g_1, \dots, g_{N_{\mathrm{GT}}}\}$ denote the
set of ground-truth connected components, and
$\mathcal{P} = \{p_1, \dots, p_{N_{\mathrm{P}}}\}$ the set of predicted
connected components.

\begin{itemize}
    \item \textbf{Lesion-wise F1 Score.} Evaluates lesion-level detection
    performance by measuring the balance between lesion-wise precision
    and recall:
    \begin{equation}
    \mathrm{Precision}_{\mathrm{les}} =
    \frac{|\mathcal{P}_{\mathrm{hit}}|}
    {|\mathcal{P}| + \varepsilon},
    \qquad
    \mathrm{Recall}_{\mathrm{les}} =
    \frac{|\mathcal{G}_{\mathrm{hit}}|}
    {|\mathcal{G}| + \varepsilon},
    \end{equation}
    \begin{equation}
    \mathrm{F1}_{\mathrm{les}} =
    \frac{
        2 \, \mathrm{Precision}_{\mathrm{les}} \, \mathrm{Recall}_{\mathrm{les}}
    }{
        \mathrm{Precision}_{\mathrm{les}} + \mathrm{Recall}_{\mathrm{les}} + \varepsilon
    },
    \end{equation}
    where $\mathcal{G}_{\mathrm{hit}} \subseteq \mathcal{G}$ and
    $\mathcal{P}_{\mathrm{hit}} \subseteq \mathcal{P}$ denote detected
    ground-truth and predicted lesions, respectively.

    \item \textbf{Small Lesion Recall.} Measures the fraction of small
    ground-truth lesions that are successfully detected:
    \begin{equation}
    \mathcal{G}_{\mathrm{small}} =
    \left\{
        g \in \mathcal{G} : |g| \leq \tau
    \right\},
    \end{equation}
    \begin{equation}
    \mathrm{Recall}_{\mathrm{small}} =
    \frac{
        \left|
        \left\{
        g \in \mathcal{G}_{\mathrm{small}} : g \in \mathcal{G}_{\mathrm{hit}}
        \right\}
        \right|
    }{
        |\mathcal{G}_{\mathrm{small}}| + \varepsilon
    }.
    \end{equation}

    \item \textbf{False Positive Volume.} Quantifies the total physical
    volume of falsely predicted lesion voxels:
    \begin{equation}
    \mathrm{FP} = P - G,
    \end{equation}
    \begin{equation}
    V_{\mathrm{FP}} =
    |\mathrm{FP}| \cdot (s_x s_y s_z),
    \end{equation}
    where $s_x, s_y, s_z$ denote voxel spacing along each axis.

    \item \textbf{False Negative Lesion Count.} Counts the number of
    ground-truth lesions that are completely missed:
    \begin{equation}
    \text{FN Lesion Count} =
    |\mathcal{G}| - |\mathcal{G}_{\mathrm{hit}}|.
    \end{equation}

    \item \textbf{False Negative Volume Fraction.} Measures the fraction
    of total lesion volume corresponding to missed lesions:
    \begin{equation}
    \mathrm{FN} = G - P,
    \end{equation}
    \begin{equation}
    V_{\mathrm{FN}}^{\mathrm{frac}}
    =
    \frac{|\mathrm{FN}|}{|G| + \varepsilon}.
    \end{equation}
\end{itemize}
